\section{Institutional background}\label{sec:background}

This section gives the institutional facts the design depends on. All of them come from public documents opened on the dates recorded in the replication files; the Thai titles are kept in the bibliography with English translations.

\subsection{The 2011 floods, insurance and relief}

The 2011 Chao Phraya floods ran from late July 2011 to January 2012 and inundated the central plain, including seven industrial estates in Ayutthaya and Pathum Thani \citep{komori2012characteristics,haraguchi2013flood}. The joint World Bank and Royal Thai Government assessment put damage and losses at THB~1.43 trillion, around 70\% in manufacturing and around 90\% borne by the private sector, and estimated that fewer than 1\% of private dwellings were insured against natural disasters \citep{worldbank2012thai}. Bank balance sheets were resilient on impact; the assessment expected non-performing loans to rise by about 0.5\% after the grace periods and restructurings that the central bank allowed.

Those grace periods are the relief rules that matter here. On 19 October 2011 the Bank of Thailand issued circular ว.~1564/2554, which let financial institutions keep the pre-flood classification of debtors who were Normal or Special Mention when the floods hit, in disaster areas declared by the Department of Disaster Prevention and Mitigation and, for indirectly affected debtors, outside them. Assisted loans were treated as general rather than troubled debt restructurings and were not reported as restructured; the attached questions and answers keep the three-month rule for upgrading Special Mention debtors, let accrual income continue, and allow collateral revaluations that fell due by 30 June 2012 to be deferred by up to a year, to 31 December 2012 at the latest \citep{bot2011circular1564}. An Emergency Decree of 26 January 2012 created a THB~300 billion soft-loan facility operated by the central bank, lending to institutions at 0.01\% with on-lending to flood-affected borrowers capped at 3\% \citep{bot2012emergency}, and a second decree created a disaster insurance promotion fund \citep{bot2012catastrophe}. A three-year moratorium at state-owned specialised financial institutions, approved by the Cabinet on 15 November 2011 for borrowers whose total principal was at most THB~500,000 on 30 June 2011, is often described as flood relief; the press report of the resolution shows it was conditioned on pre-flood debt status, not on flooding, and I do not use it as a flood instrument \citep{mgr2011cabinetmoratorium}.\footnote{The resolution text itself could not be opened. The World Bank assessment describes the moratorium as flood relief \citep[Table~21]{worldbank2012thai}; the press report of the Cabinet resolution gives eligibility as NPL, restructured or (for the agricultural bank) delinquent status with total principal of at most THB~500,000 at 30 June 2011, with no flood condition.}

The 2011 template became a standing arrangement in 2024. Since 30 August 2024 circular ว.~601/2567, framed as a request for cooperation, has let lenders keep the pre-disaster classification of assisted debtors who are not yet credit-impaired, for up to twelve months from a disaster-zone declaration or from a yellow, orange or red disaster-alert level \citep{bot2024circular601,bot2024relief}. After the 2025 floods, circular ว.~8158/2568 of 15 December 2025 added, for Level-4 (``extremely severe'') zones, a principal moratorium and interest waiver of up to twelve months for retail and SME debtors who were not non-performing on 2 December 2025, with opt-in by 31 January 2026; days past due may be held at the pre-disaster count, home-repair loans are excluded from the loan-to-value calculation and risk-weighted at 35\% \citep{bot2025level4}. Three features of these rules recur in the identification strategy of Section~\ref{sec:credit}: eligibility turns on pre-flood classification status; the relief window is keyed to the date and place of an administrative declaration rather than to flooding; and the rules freeze the observed default of eligible borrowers for a known period.

\subsection{Official land values}

The Treasury Department sets an official value for every titled land parcel under the Property Valuation for State Benefit Act B.E.~2562 (2019) \citep{th2019valuationact}. Values are set in rounds of four years, published at least thirty days before they take effect, and a round may be extended by one year (s.~19). The method is laid down in a 2020 Ministerial Regulation: provincial committees value land by market comparison from registered transactions and other market evidence, assign a value per square wah to each land unit, and adjust for depth and shape by fixed factors, with values rounded to the schedule \citep{th2020appraisalmr}. The regulation contains no flood or hazard rule; its one explicit trigger for a within-round revision is a divergence of 15\% or more between official and market values (ข้อ~9). The values are the base for the 2\% transfer fee and for other fees and taxes collected at registration \citep{dol2026fees}. The current round covers 1 January 2023 to 31 December 2026; the parcel-level schedule is published as open data \citep{treasury2026opendata}.

Two points follow. Because the method is market comparison, the official values should capitalise whatever registered transaction prices capitalise, with a lag set by the round calendar and by the committees' sampling of comparables. The one direct piece of evidence on how the Department responds to a flood points to a lag of at least a round: after the 2011 floods it postponed the new schedule by six months, reported that its survey found no fall in prices in flooded areas, and brought the schedule compiled in late 2011 into force unchanged, with an average rise of 21.34\% \citep{mgr2012appraisal}. And because the calendar is fixed in advance, floods that occur between rounds, such as the 2024 northern and 2025 floods, create a difference-in-differences design: parcels flooded between two rounds against nearby unflooded parcels, with the value change between rounds as the outcome.

\subsection{Bank collateral values, classification and capital}

The Bank of Thailand's collateral policy of 8 May 2019 governs the values that institutions use in provisioning \citep{bot2019collateral}. Credit-impaired (stage~3) real-estate collateral must be revalued at least every three years; for performing and under-performing collateral the frequency is at the institution's discretion, and automated valuation models are allowed subject to validation. Methods follow the valuation standards endorsed by the securities regulator. The policy makes no reference to the Treasury Department's official values; a full-text search of the document found none. The link between official values and bank collateral values is therefore not a rule but, at most, a shared reliance on the same transaction evidence; Section~\ref{sec:framework} treats it as an assumption.

Classification and provisioning moved to the three TFRS~9 stages, performing, under-performing and credit-impaired, for financial statements beginning on or after 1 January 2020 \citep{bot2018tfrs9}. The 2011 floods therefore happened under the earlier five-class system, and the 2024--25 floods under expected-credit-loss staging with a classification freeze. Capital thresholds are a 4.5\% common equity tier~1 minimum plus a 2.5\% conservation buffer \citep{bot2012capital}, with a further 1\% for domestic systemically important banks \citep{bot2017dsib}; these give the 7\% and 8\% thresholds used in Section~\ref{sec:stress}.

\subsection{Climate supervision}

The Bank of Thailand's Standard Practice of 15 February 2023 asks financial institutions to integrate environmental and climate considerations into their business; flood is named as an acute physical risk, scenario analysis appears as an example of forward-looking assessment without a prescribed scenario, and incorporating environmental risk into collateral valuation is given as an example control for high-risk customers \citep{BOT_StandardPractice_page}. A pilot climate stress test with systemically important and other ready banks, ``starting with physical risk'', is listed under the central bank's 2024 progress without published scenarios or results \citep{bot2025sustainability2567}. A World Bank background note for Thailand's Country Climate and Development Report, prepared with the central bank, maps corporate lending by district flood risk and finds 41\% of the sample banks' corporate lending in districts with high riverine flood risk and 7\% in very-high-risk districts, with a 2011-type flood treated as a 1-in-50-year event \citep{worldbank2026ccdrfin,dolk2025assessing}. These documents set the supervisory template that Section~\ref{sec:stress} formalises: hazard footprint, collateral damage, loan-to-value, loss given default, capital.
